\scanpage{112}
\mspage{6}
% archive foliation: 111

Montrons comment cela permet de calculer $z$, \unsure{i.e}
$z_0=\wp(z)$ et $z'=\wp'(z)$ (reliés par
\[
(46)\qquad z'^2=z_0^3+a_1z_0^2+a_2z_0+a_3),
\]
quand on connaît $x$ et $y$, exprimés à l'aide de
\deleted{$x_0=\wp(x),x'=\wp'(x)$}
\[
\left\{\begin{aligned}
x_0&=\wp(x),\ x'=\wp'(x)\\
y_0&=\wp(y),\ y'=\wp'(y)
\end{aligned}\right.
\]
(reliés par $x'^2=x_0^3+a_1x_0^2+a_2x_0+a_3$, $y'^2=y_0^3+a_1y_0^2+a_2y_0+a_3$).

L'équation (\unsure{41|42}) s'écrit
\[
(*)\qquad z'(x_0-y_0)+z_0(y'-x')+(x'y_0-y'x_0)=0
\]
i.e.
\[
(47)\qquad z'=\lambda z_0+\mu,\quad\text{avec }\lambda=\frac{x'-y'}{x_0-y_0},\quad
\mu=\frac{y'x_0-x'y_0}{x_0-y_0}.
\]
\deleted{Donc} (NB on a $x_0\ne y_0$ i.e $x_0-y_0\ne0$).
\deleted{De} plus $z'$ satisfait (46), qui équivaut à
\deleted{l'équ} (compte tenu de (47), qui implique
$z'^2=\lambda^2z_0^2+2\lambda\mu z_0+\mu^2$) : l'équation
\[
(48)\qquad \unsure{z^3|z_0^3}+(a_1-\lambda^2)z_0^2+(a_2-2\lambda\mu)z_0+(a_3-\mu^2)=0.
\]
Donc $(z_0,z')$ est lié par (47), (48), a priori on trouve 3 solutions
\inserted{en $z_0$} de (48), à qui correspondent trois solutions $(z_0,z')$ via (47).
Il faut trouver une solution unique, en se rappelant qu'on doit avoir
\[
(49)\qquad z_0\ne x_0,\quad z_0\ne y_0,\qquad\deleted{\text{\unsure{alors}}}
\]
